\section{Results}

\subsection{Model Comparison}

Table~\ref{tab:results} reports suggestion-class F1 for each
model on both datasets. On the SemEval test set, BERT stage-1
achieves the highest suggestion-class F1 (0.716), followed by
TF-IDF~+~LR (0.596) and the regex baseline (0.386). Each tier
improves over the previous, confirming that learned
representations outperform hand-crafted patterns for in-domain
suggestion detection.

\begin{table}[t]
\centering\footnotesize
\caption{Suggestion-class metrics. P = precision, R = recall.
  $\Delta$ = MBS F1 $-$ SemEval F1.}
\label{tab:results}
\begin{tabular}{l@{\hskip 4pt}ccc@{\hskip 6pt}cccc}
\toprule
& \multicolumn{3}{c}{\textbf{SemEval}}
& \multicolumn{3}{c}{\textbf{MBS}} & \\
\cmidrule(lr){2-4} \cmidrule(lr){5-7}
\textbf{Model} & P & R & F1 & P & R & F1
  & $\Delta$ \\
\midrule
Regex      & .405 & .368 & .386 & .421 & .533 & .471
  & $+$.09 \\
TF-IDF+LR  & .486 & .770 & .596 & .333 & .133 & .190
  & $-$.41 \\
BERT s-1   & .624 & .839 & .716 & .500 & .200 & .286
  & $-$.43 \\
BERT s-2   & .513 & .908 & .656 & .563 & .600 & .581
  & $-$.08 \\
\bottomrule
\end{tabular}
\end{table}

On the MBS test split, BERT stage-2 achieves the best
suggestion-class F1 (0.581), approximately doubling the stage-1
result (0.286). The two-stage fine-tuning improves both precision (0.500
$\rightarrow$ 0.563) and recall (0.200 $\rightarrow$ 0.600) on
MBS, indicating that domain-specific training helps the model
recognise hotel-style suggestion language.
Figure~\ref{fig:cm-s2} shows the confusion matrix for BERT
stage-2 on the MBS test split.

\begin{figure}[t]
\centering
\includegraphics[width=0.75\columnwidth]{cm_bert_s2_mbs.png}
\caption{Confusion matrix for BERT stage-2 on the MBS test
  split (80 sentences).}
\label{fig:cm-s2}
\end{figure}

\subsection{Domain Transfer Analysis}

Cross-domain performance reveals a striking pattern: the two
learned models suffer large drops when moving from SemEval to MBS
(TF-IDF~+~LR: $-$0.406; BERT stage-1: $-$0.430), while the regex
baseline actually improves slightly ($+$0.085). This suggests that
the learned models overfit to SemEval-specific vocabulary and
sentence structures (e.g.\ software feature requests), whereas
simple modal and conditional patterns transfer across domains.
The differing suggestion base rates between the SemEval test set
(10.4\%) and MBS test split (18.8\%) may also partly explain the
regex improvement, as a higher base rate favours
pattern-matching approaches.

However, the regex baseline's cross-domain robustness comes at
the cost of low absolute performance (suggestion-class F1 of
0.471 on MBS). The regex approach generates 11 false positives on
the 80-sentence MBS test split, primarily from positive modal
phrases (e.g.\ ``We would recommend this hotel to anyone'')
that match suggestion patterns but express praise.

BERT stage-2 demonstrates that even a small amount of in-domain
labelled data (320 sentences) can substantially recover
cross-domain performance. Compared to BERT stage-1 on MBS,
stage-2 reduces false negatives from 12 to 6 and improves recall
from 0.200 to 0.600, at the cost of a modest increase in false
positives (3 $\rightarrow$ 7). On the SemEval test set, stage-2
suggestion-class F1 drops from 0.716 to 0.656, with precision
falling (0.624 $\rightarrow$ 0.513) while recall rises (0.839
$\rightarrow$ 0.908)---moderate forgetting from hotel
fine-tuning. Figure~\ref{fig:model-comparison} summarises
suggestion-class F1 across models and datasets.

\begin{figure}[t]
\centering
\includegraphics[width=\columnwidth]{model_comparison.png}
\caption{Suggestion-class F1 by model and dataset.}
\label{fig:model-comparison}
\end{figure}

\subsection{Suggestion Themes (BERTopic)}

BERTopic analysis of 4,179 predicted suggestions from BERT
stage-2 identified 56 topics, with topic~0 containing 2,278
sentences and the remaining 55 topics ranging from 5 to 56
sentences each.
Topic~0 functions as a catch-all for sentences lacking
distinctive sub-themes; the finer-grained topics (1--55) capture
specific operational areas such as minibar pricing, elevator
access, and pool amenities.

\begin{table}[t]
\centering\small
\caption{Aspect distribution of predicted suggestions.}
\label{tab:aspects}
\begin{tabular}{lrr}
\toprule
\textbf{Aspect} & \textbf{Count} & \textbf{\%} \\
\midrule
Other            & 1,299 & 31.1 \\
Room             &   990 & 23.7 \\
Facilities       &   638 & 15.3 \\
Food \& beverage &   490 & 11.7 \\
Value            &   309 &  7.4 \\
Service          &   234 &  5.6 \\
Location         &   219 &  5.2 \\
\bottomrule
\end{tabular}
\end{table}

Room-related suggestions dominate (23.7\%), followed by
facilities (15.3\%) and food and beverage (11.7\%). The ``other''
category (31.1\%) captures suggestions that do not match
predefined aspect keywords. Room amenities and physical
facilities together account for 39\% of all detected
suggestions, providing a clear priority for hotel management.
